% STATUS: VERIFIED / AUDITED
% VERSION: P02 v1.1; 2026-09-27; A02 E01 corrected
% AUDIT FILE: audits/A02_synchronous_subspace_and_equilibria.tex
% AUDIT EVIDENCE: audits/A02_synchronous_subspace_and_equilibria.pdf
% AUDIT SOURCE AVAILABILITY: Only the PDF was supplied; the TeX source is absent.
% ORIGINAL AUDIT VERDICTS: CORE MINOR REVISION; EXTENDED VERIFIED
% CORE VERDICT: VERIFIED AFTER MINOR REVISION
% EXTENDED VERDICT: VERIFIED
% CLOSURE: A02 E01 corrected; boundary limits explicitly name (s_-,s_+).
% BODY SHA-256: 8a061ba81c7ffd2d05b188b9fca4ec1b94bc132147f7f664c8523f0a5f32560c
% HASH SCOPE: UTF-8 bytes from the first \section{P02: line through end of file.
% STATUS NOTE: Historical DRAFT/audit-pending body wording is preserved;
% this header and the project manifest govern current status.
% DEPENDS ON: P01 CORE VERIFIED / AUDITED; canonical model and notation
% MODEL: canonical demand-modulated model
% LITERATURE DEPENDENCIES: none
% NUMERICAL DEPENDENCIES: none
% NEXT DEPENDENCIES: P03, P04, P09, P10, P11
% SCOPE: synchronous invariance, scalar restriction, complete canonical root
% classification, domain filtering, and algebraic branch loci only.
% CLAIMS: C025, C026, C028--C033, C146; no later-module claim is promoted.
% CHECKS: exact algebra and representative/equality cases; no simulations.

\section{P02: Synchrony and synchronous equilibria}
\label{sec:P02}

\noindent\textbf{Status: DRAFT; independent audit pending.}
The standing assumptions are the canonical A1--A3. In particular,
$n\ge5$, $1\le q\le n-1$, $a,k,c,g,h>0$, and the open extension
domain is $\Domain=(-1,\infty)^n$. The nonnegative model state space
is $\Orth$. P01 v1.1 is VERIFIED / AUDITED: its required wording
correction E01 is already incorporated. Only its local uniqueness
and, for nonnegative initial data, global continuation are needed here.
The calculations below start from \eqref{eq:canonical-dm}.
No assertion about nonsynchronous equilibria, equilibrium stability,
or bifurcation type is made.

\subsection{Invariant synchronous subspace}
\label{slot:P02-1}

\begin{proposition}[Invariance of the synchronous subspace]
\label{prop:P02-invariance}
The set
\[
 \Sync\cap\Domain=\{s\one:s\in(-1,\infty)\}
\]
is invariant for the canonical vector field for the entire maximal
solution interval in $\Domain$. Its scalar restriction is
\begin{equation}
\label{eq:P02-scalar-original}
 \dot s=f(s):=as-ks^2+\frac{cs}{1+s}
                     -\frac{gs}{1+s}-hs^2,
 \qquad s\in(-1,\infty).
\end{equation}
This conclusion needs no coprimality or distinct-offset assumption.
Indeed, the substitution and local invariance remain valid for any
integer $n\ge1$, any residue $q$ modulo $n$, and any finite real
coefficients. The canonical restrictions are retained for the
equilibrium classification below.
\end{proposition}

\begin{proof}
Put $x=s\one$ with $s>-1$. For every $i\in\Zn$, each of
$x_i,x_{i-1},x_{i+q},x_{i-q}$ equals $s$. Direct substitution gives
\[
 F_i(s\one)=as-ks^2+\frac{cs}{1+s}
                        -\frac{gs}{1+s}-hs^2=f(s).
\]
Every component is identical, so $F(s\one)=f(s)\one$.
Coincident offsets retain all the displayed terms; no distinctness
of neighbors is used. Nor is the size or connectedness of a
step-$q$ cycle used.

The scalar field is smooth on $(-1,\infty)$. A scalar solution
$s(t)$ lifts to the vector solution $s(t)\one$. Uniqueness from
Proposition~\ref{prop:P01-local} makes this lift agree with the
vector solution having initial value $s(0)\one$ wherever both exist.
If the scalar maximal interval ended earlier while the vector
solution continued in $\Domain$, the latter would supply a limit
$s_*\one\in\Domain$ at that time. Local scalar continuation at
$s_*>-1$ would extend the lift, a contradiction. Conversely, a longer
scalar interval would extend the vector solution. Thus the maximal
intervals coincide and synchrony persists throughout. The same local
uniqueness argument applies on either side of an initial time.
The scalar domain inherited from $\Domain$ is exactly $(-1,\infty)$.
\end{proof}

\subsection{Scalar synchronous dynamics}
\label{slot:P02-2}
Use the reserved abbreviations and one local auxiliary abbreviation
\begin{equation}
\label{eq:P02-abbreviations}
 A=k+h>0,\qquad \mu=a+c-g,\qquad b=a-A.
\end{equation}
Here $b$ is used only within P02; it is not the archived comparator's
coefficient. In particular $b>-A$ because $a>0$. The exact restriction is
\begin{equation}
\label{eq:P02-scalar}
 f(s)=s\left(a-As+\frac{c-g}{1+s}\right)
     =\frac{s}{1+s}\bigl(\mu+bs-As^2\bigr),\qquad s>-1.
\end{equation}
For $s(0)\ge0$, P01 and Proposition~\ref{prop:P02-invariance}
give a global nonnegative scalar solution. For $-1<s(0)<0$, only
the maximal interval within this open domain is asserted. Even when
$c=g$ cancels the two rational terms along synchrony, this does not
enlarge the canonical domain or admit $s=-1$.

\subsection{The origin and the nonzero equilibrium equation}

\begin{proposition}[Equilibrium equation with domain restrictions]
\label{prop:P02-equation}
The origin $0\one$ is an equilibrium for every canonical parameter
tuple. Every other synchronous equilibrium in $\Domain$ is exactly
$s\one$, where $s\ne0$, $s>-1$, and
\begin{align}
 0=Q(s)&:=As^2-bs-\mu \label{eq:P02-polynomial}\\
       &=(k+h)s^2+(k+h-a)s-(a+c-g).\nonumber
\end{align}
Define the discriminant
\begin{align}
 D&:=b^2+4A\mu \label{eq:P02-discriminant}\\
  &=(a-A)^2+4A(a+c-g)
   =(a+A)^2+4A(c-g).\nonumber
\end{align}
The complete list of real polynomial roots, when $D\ge0$, is
\begin{equation}
\label{eq:P02-roots}
 s_\pm=\frac{b\pm\sqrt D}{2A}
 =\frac{a-(k+h)\pm
 \sqrt{(a+k+h)^2+4(k+h)(c-g)}}{2(k+h)}.
\end{equation}
They are distinct for $D>0$ and coincide for $D=0$. When $D<0$
there are no real roots of $Q$.
\end{proposition}

\begin{proof}
Directly in \eqref{eq:P02-scalar-original}, all five terms vanish
at $s=0$. Thus the origin is always an equilibrium, before any
division is performed. Now take $s\ne0$ with $s>-1$.
Dividing \eqref{eq:P02-scalar} by $s$ and multiplying by
the strictly positive number $1+s$ gives
\[
 (a-As)(1+s)+c-g=\mu+(a-A)s-As^2=0.
\]
Both operations are reversible on this stated set, giving
\eqref{eq:P02-polynomial}. Completing the square yields
\begin{equation}
\label{eq:P02-square}
 Q(s)=A\left(s-\frac b{2A}\right)^2-\frac D{4A}.
\end{equation}
Expansion of $b^2+4A\mu$ proves both discriminant forms and the
quadratic formula gives \eqref{eq:P02-roots}. Finally,
\begin{equation}
\label{eq:P02-minusone}
 Q(-1)=A+b-\mu=g-c.
\end{equation}
Thus $s=-1$ is a polynomial root precisely when $g=c$, but it is
never a point of the canonical domain. Multiplication by $1+s$
does not authorize retaining that root. A root $s=0$ of $Q$, when
present, is the already counted origin rather than another equilibrium.
\end{proof}

\subsection{Complete classification in the canonical parameter regime}
\label{slot:P02-3}

In this subsection, a \emph{mathematical root} means a real root of
$Q$, including a zero or an inadmissible root. A \emph{DM-domain root}
additionally satisfies $s>-1$. A \emph{nonnegative model equilibrium}
has $s\ge0$, and a \emph{positive model equilibrium} has $s>0$.
The ever-present origin is added once to the admissible nonzero roots.
All counts below count distinct values, not algebraic multiplicity.

\begin{theorem}[Complete synchronous-equilibrium classification]
\label{thm:P02-complete}
Under the canonical assumptions, the synchronous equilibrium set is
\begin{equation}
\label{eq:P02-equilibrium-set}
 \{0\one\}\ \cup\
 \{s\one:Q(s)=0,\ s>-1,\ s\ne0\}.
\end{equation}
The following eleven cases are exhaustive. In the table, $r$ counts
real nonzero polynomial roots, $d$ counts those with $s>-1$, $N_0$
counts all nonnegative synchronous equilibria (including the origin),
and $N_+$ counts the strictly positive ones. The total number of
synchronous equilibria in $\Domain$ is $1+d$.

\begin{center}
\small
\renewcommand{\arraystretch}{1.18}
\begin{tabular}{@{}c p{0.28\textwidth} p{0.35\textwidth} c c c c@{}}
\toprule
Case & Parameter conditions & Real roots of $Q$ and location
 & $r$ & $d$ & $N_0$ & $N_+$\\
\midrule
I & $\mu>0,\ g<c$ & $s_-<-1<0<s_+$ & 2&1&2&1\\
II & $\mu>0,\ g=c$ & $s_-=-1,\ s_+=a/A>0$ & 2&1&2&1\\
III & $\mu>0,\ g>c$ & $-1<s_-<0<s_+$ & 2&2&2&1\\
IV & $\mu=0,\ b>0$ & $s_-=0,\ s_+=b/A>0$ & 1&1&2&1\\
V & $\mu=0,\ b=0$ & $s_-=s_+=0$ (double) & 0&0&1&0\\
VI & $\mu=0,\ b<0$ & $-1<s_-=b/A<0=s_+$ & 1&1&1&0\\
VII & $\mu<0,\ D<0$ & No real roots; any sign of $b$ & 0&0&1&0\\
VIII & $\mu<0,\ D>0,\ b>0$ & $0<s_-<s_+$ & 2&2&3&2\\
IX & $\mu<0,\ D=0,\ b>0$ & $s_-=s_+=b/(2A)>0$ & 1&1&2&1\\
X & $\mu<0,\ D>0,\ b<0$ & $-1<s_-<s_+<0$ & 2&2&1&0\\
XI & $\mu<0,\ D=0,\ b<0$ & $-\tfrac12<s_-=s_+=b/(2A)<0$ & 1&1&1&0\\
\bottomrule
\end{tabular}
\end{center}
When $\mu<0$ and $b=0$, necessarily $D=4A\mu<0$, so case VII
applies. Each appearance of $s_\pm$ refers to \eqref{eq:P02-roots}.
The only repeated polynomial roots are in cases V, IX and XI.

As zeros of the original smooth scalar field $f$ on $(-1,\infty)$,
all nonzero equilibria are simple except the double roots in IX and XI.
The origin is simple if $\mu\ne0$, double if $\mu=0$, $b\ne0$,
and triple if $\mu=b=0$. The algebraic continuation of the nonzero
equilibrium curve meets the origin exactly at $\mu=0$: through $s_-$
in case IV, through $s_+$ in case VI, and through both coincident
root labels in case V. No negative root represents a state in $\Orth$.
\end{theorem}

\begin{proof}
Proposition~\ref{prop:P02-equation} proves completeness of
\eqref{eq:P02-equilibrium-set}. The Vieta relations are
\begin{equation}
\label{eq:P02-vieta}
 s_-+s_+=\frac bA,\qquad s_-s_+=-\frac\mu A.
\end{equation}
We first establish the domain test for every real pair of roots.
Because $a>0$ and $A>0$,
\[
 s_+\ge\frac b{2A}=\frac a{2A}-\frac12>-\frac12>-1.
\]
For the smaller root,
\begin{equation}
\label{eq:P02-domain-test}
 s_->-1
 \ \Longleftrightarrow\ \sqrt D<a+A
 \ \Longleftrightarrow\ D<(a+A)^2
 \ \Longleftrightarrow\ g>c.
\end{equation}
The middle equivalence uses $a+A>0$ and $D\ge0$. Equality gives
$s_-=-1$ exactly when $g=c$; the reverse inequality gives $s_-<-1$
exactly when $g<c$. These are polynomial-root statements on the
real-discriminant set, not assertions that such roots are admissible.

If $\mu>0$, then $D>b^2$, so $s_-<0<s_+$ and both roots are
nonzero and simple. The domain test gives I--III. When $g=c$,
$D=(a+A)^2$ and the stated roots $-1,a/A$ follow directly.

If $\mu=0$, then $Q(s)=s(As-b)$. If $b\ne0$, its distinct roots
are $0$ and $b/A$. Since $b/A=a/A-1>-1$, the latter is always
domain-admissible and has the sign of $b$. Their ordering gives
IV and VI. If $b=0$, $Q(s)=As^2$ gives V. The origin is counted
only once in the equilibrium set in all three cases.

If $\mu<0$, real roots exist precisely when $D\ge0$. For $D<0$
we obtain VII. For $D\ge0$, the product in \eqref{eq:P02-vieta}
is strictly positive and $b\ne0$ (otherwise $D<0$). Thus both
roots are nonzero with the sign of their sum, namely the sign of $b$.
They are distinct for $D>0$ and double for $D=0$. Also
$g>a+c>c$, so both are admissible by \eqref{eq:P02-domain-test}.
This gives VIII--XI, including the sharper negative repeated-root
bound from $b>-A$. These arguments yield every count in the table.

On $s>-1$, the identity $f(s)=-sQ(s)/(1+s)$ shows that the
nonzero zeros have the same multiplicity as their polynomial roots.
At the origin, use \eqref{eq:P02-scalar} without dividing:
the numerator is $\mu s+bs^2-As^3$ and $1+s$ does not vanish.
This gives the stated orders one, two and three. Finally $Q(0)=-\mu$,
so the nonzero equation's algebraic continuation contains $s=0$
exactly when $\mu=0$. The formulas at that locus give its root
labels, and the local branch description below confirms the approach
by nonzero roots. The counting and multiplicity assertions are purely
algebraic and carry no stability conclusion.
\end{proof}

\subsection{Positive synchronous equilibria: necessary and sufficient conditions}

\begin{theorem}[Exact positive-equilibrium alternatives]
\label{thm:P02-positive}
In the canonical parameter regime the following four alternatives
partition parameter space:
\begin{enumerate}[label=(\roman*),leftmargin=*]
\item There is no positive synchronous equilibrium if and only if
\begin{equation}
\label{eq:P02-none}
 [\,\mu=0,\ b\le0\,]
 \quad\hbox{or}\quad
 [\,\mu<0,\ (b\le0\ \hbox{or}\ D<0)\,].
\end{equation}
\item There is exactly one positive synchronous equilibrium and it is
simple if and only if $\mu>0$, or $\mu=0$ and $b>0$.
It is $s_+\one$; in the second case $s_+=b/A$.
\item There are exactly two distinct positive synchronous equilibria
if and only if $\mu<0$, $b>0$, and $D>0$. They are
$s_-\one,s_+\one$, both simple.
\item There is exactly one positive synchronous equilibrium and it
is repeated if and only if $\mu<0$, $b>0$, and $D=0$.
It is a double root, $s_-\one=s_+\one=(b/(2A))\one$.
\end{enumerate}
Consequently, exactly one distinct positive equilibrium (without a
simplicity qualification) occurs if and only if (ii) or (iv) holds.
There are never more than two positive synchronous equilibria.
\end{theorem}

\begin{proof}
The cases of Theorem~\ref{thm:P02-complete} prove every implication,
including the necessity and all equalities. One can also reconstruct
the positive-root count directly from
$H(s)=\mu+bs-As^2$ for $s>0$. If $b\le0$, then
$H'(s)=b-2As<0$ on this interval: it has one positive zero exactly
when $\mu>0$. If $b>0$, its maximum on $s\ge0$ is at
$s=b/(2A)>0$ and has value $D/(4A)$. For $\mu<0$, that maximum
is positive, zero, or negative precisely in alternatives (iii),
(iv), or (i). For $\mu=0$, the zero at the endpoint is excluded
from the positive count and the other root is $b/A>0$. For
$\mu>0$, the opposite-sign Vieta product leaves one positive root.
This gives an independent sign-and-vertex justification of the list.
\end{proof}

\subsection{Algebraic branch loci and degeneracies}

\begin{proposition}[Branch-intersection and repeated-root loci]
\label{prop:P02-loci}
For fixed positive $a,k,c,h$, define
\begin{equation}
\label{eq:P02-gf}
 g_f:=a+c+\frac{(a-A)^2}{4A}
      =c+\frac{(a+A)^2}{4A},\qquad
 s_f:=\frac{a-A}{2A}.
\end{equation}
The following statements concern only the scalar equilibrium equation.
\begin{enumerate}[label=(\roman*),leftmargin=*]
\item The \emph{branch-intersection locus} is $\mu=0$, or $g=a+c$.
It is the set where the algebraic continuation of the nonzero
equilibrium equation meets the origin branch.
\item The \emph{repeated-root locus} of $Q$ is $D=0$, or $g=g_f$,
and its root is $s_f$. Indeed,
\begin{equation}
\label{eq:P02-discriminant-factor}
 D=4A(g_f-g),\qquad
 g_f-(a+c)=\frac{(a-A)^2}{4A}\ge0.
\end{equation}
The root is positive for $a>A$, negative for $0<a<A$, and zero
for $a=A$. It always lies in $(-1/2,\infty)$ and hence in the
DM domain. It is a nonzero double equilibrium only when $a\ne A$.
\item The two parameter loci meet exactly at the \emph{algebraic
degeneracy} $a=A$, $g=a+c$. There $Q(s)=As^2$ and
\begin{equation}
\label{eq:P02-degenerate-field}
 f(s)=-\frac{As^3}{1+s}.
\end{equation}
The origin is the sole synchronous equilibrium in the domain at
that parameter tuple; it is a triple zero of $f$ and a double root
of $Q$, not an additional nonzero equilibrium.
\end{enumerate}
\end{proposition}

\begin{proof}
The nonzero equation can be written, without approximation, as
\begin{equation}
\label{eq:P02-branch-graph}
 \mu=As^2-bs,\qquad
 g=a+c+bs-As^2
   =g_f-A(s-s_f)^2.
\end{equation}
This is its algebraic continuation, to be filtered by $s>-1$,
$g>0$, and $s\ne0$ when counting nonzero model equilibria.
At $s=0$ it gives exactly $g=a+c$. Completing the square or using
\eqref{eq:P02-square} gives $D=0$, $g=g_f$, and $s=s_f$ as the
unique repeated-root condition. The difference of the two parameter
values in \eqref{eq:P02-discriminant-factor} vanishes exactly at
$a=A$. Since $a>0$, $s_f=a/(2A)-1/2>-1/2$. At the common
locus, substitution into \eqref{eq:P02-scalar} proves
\eqref{eq:P02-degenerate-field} and all the remaining assertions.
\end{proof}

For additional local algebra at the origin, fix $a,A,c$ and vary $g$
(equivalently $\mu$). If $b\ne0$, the root that tends to zero is
\begin{equation}
\label{eq:P02-local-branch}
 s_0(\mu)=\frac{b-\operatorname{sgn}(b)\sqrt{b^2+4A\mu}}{2A}
        =-\frac{\mu}{b}+\frac{A\mu^2}{b^3}+O(\mu^3).
\end{equation}
This follows by expansion of the explicit square root near $\mu=0$;
$g=a+c-\mu$ remains positive in a sufficiently small neighborhood.
Thus nonzero domain-admissible roots approach the origin from the
neighboring parameter regions. At the intersection itself,
\begin{equation}
\label{eq:P02-quadratic-coefficient}
 f(s)\big|_{\mu=0}=\frac{s^2(b-As)}{1+s}
                     =bs^2-as^3+O(s^4).
\end{equation}
The coefficient of the usual quadratic scalar term is $b=a-A$.
It vanishes when $a=A$; this is precisely the degeneracy above.

More explicitly, for $a=A$ the nonzero polynomial equation is
$As^2=\mu$. For $\mu>0$ its roots are
\begin{equation}
\label{eq:P02-square-root-branches}
 s_\pm=\pm\sqrt{\mu/A}.
\end{equation}
The negative one is admissible exactly when $0<\mu<A$ (equivalently
$c<g<a+c$); it is $-1$ when $\mu=A$, and below $-1$ when $\mu>A$.
At $\mu=0$ both labels give the origin; for $\mu<0$ no real
nonzero roots exist. These are algebraic branch descriptions only.
Any bifurcation theorem requires the separate work assigned to P09
and P10; none is asserted here.

\subsection{Canonical parameter corollary}

\begin{corollary}[Positive branches in the original canonical parameters]
\label{cor:P02-canonical}
Let $a,k,c,g,h>0$, $A=k+h$, and let $g_f$ be as in
\eqref{eq:P02-gf}. If $a\le k+h$, there is one simple positive
synchronous equilibrium exactly for $0<g<a+c$, and none for
$g\ge a+c$. If $a>k+h$, the exact alternatives are
\begin{center}
\begin{tabular}{@{}ll@{}}
\toprule
Parameter range & Positive synchronous equilibria\\
\midrule
$0<g\le a+c$ & One simple equilibrium, $s_+\one$\\
$a+c<g<g_f$ & Two distinct simple equilibria, $s_-\one,s_+\one$\\
$g=g_f$ & One double equilibrium, $s_f\one$\\
$g>g_f$ & None\\
\bottomrule
\end{tabular}
\end{center}
In particular, $a+c<g$ does not by itself preclude positive
synchrony: two distinct positive equilibria exist if and only if
\begin{equation}
\label{eq:P02-two-positive-original}
 a>k+h,\qquad
 a+c<g<a+c+\frac{(a-k-h)^2}{4(k+h)}.
\end{equation}
The upper equality gives one repeated positive equilibrium when
$a>k+h$; the lower equality gives one simple positive equilibrium
and the origin. The origin exists in every row, independently of
the displayed positive count. There is no dependence on $n$, $q$,
$\gcd(n,q)$, or distinctness of the offsets.
\end{corollary}

\begin{proof}
The sign of $b$ is the sign of $a-(k+h)$; the sign of $\mu$ is
the sign of $a+c-g$; and $D=4A(g_f-g)$ with $A>0$.
Substitute these identities into Theorem~\ref{thm:P02-positive}.
\end{proof}

\subsection{Weaker assumptions and limits of the classification}
The synchronous substitution and the scalar identity require only
homogeneity, finite real coefficients, and nonvanishing denominators.
They therefore apply under P01's weaker assumptions as well.
For $A>0$, the polynomial formula, discriminant, Vieta relations,
and the positive-root alternatives of Theorem~\ref{thm:P02-positive}
remain valid for real $a,c,g$: their positive-root proofs use only
$A>0$. This observation is not a generalization of the canonical
negative-root table. In particular, the domain test
\eqref{eq:P02-domain-test} uses $a+A>0$, and the bound $s_f>-1/2$
uses $a>0$. P01 allows real $a$ and does not supply these inequalities.
The complete domain-filtered classification in this module is
intentionally restricted to the canonical positive parameter regime.
No classification for $A=0$ or $A<0$ is claimed.

\subsection{Checks and possible failure points}
\label{subsec:P02-checks}

\paragraph{Root substitution and discriminant reconstruction.}
For $D\ge0$, put $d=\sqrt D$ and
$r=(b+\varepsilon d)/(2A)$, $\varepsilon\in\{-1,1\}$.
Direct expansion gives
\[
 \mu+br-Ar^2
 =\mu+\frac{2b(b+\varepsilon d)-(b+\varepsilon d)^2}{4A}
 =\mu+\frac{b^2-d^2}{4A}=0.
\]
For $r\ne-1$, substitution in the original scalar expression gives
\[
 ar-kr^2+\frac{cr}{1+r}-\frac{gr}{1+r}-hr^2
   =\frac{r}{1+r}(\mu+br-Ar^2)=0.
\]
Only $r>-1$ is retained as a DM equilibrium. If $r=-1$, the two
original rational terms are undefined even when $c=g$, so no
substitution or canceled extension admits that point.
At $\mu=0$, the specialized root $r=b/A$ also gives
$br-Ar^2=0$. At $D=0$, $r=s_f=b/(2A)$ gives
$\mu+br-Ar^2=\mu+b^2/(4A)=0$. At $g=c$, the retained root
$a/A$ gives $ar-Ar^2=0$. At $b=0$, the square-root branches give
$\mu-Ar^2=0$. Finally, direct substitution at $r=0$ gives zero
for every parameter tuple. These checks include every exact root
formula displayed above, with $s_0$ in \eqref{eq:P02-local-branch}
being one of $s_\pm$; its truncated expansion is not used as a root.

Independently expanding the polynomial discriminant gives
\[
 (A-a)^2-4A[-(a+c-g)]
 =a^2+2aA+A^2+4A(c-g)
 =(a+A)^2+4A(c-g).
\]
Also the sum of the two explicit roots is $b/A$, and their product
is $(b^2-D)/(4A^2)=-\mu/A$, confirming both Vieta relations.

\paragraph{Exact representative cells and equality boundaries.}
The following rational tuples realize every row of the classification.
Set $k=h=1/2$, hence $A=1$, and $c=1$ throughout.
All original coefficients are strictly positive. Roots are obtained
by the formulas already proved; the counts refer to the preceding table.
\begin{center}
\small
\renewcommand{\arraystretch}{1.12}
\begin{tabular}{@{}c c c c c c@{}}
\toprule
Case & $a$ & $g$ & $\mu$ & $D$ & $(r,d,N_0,N_+)$\\
\midrule
I & $2$ & $1/2$ & $5/2$ & $11$ & $(2,1,2,1)$\\
II & $2$ & $1$ & $2$ & $9$ & $(2,1,2,1)$\\
III & $2$ & $2$ & $1$ & $5$ & $(2,2,2,1)$\\
IV & $2$ & $3$ & $0$ & $1$ & $(1,1,2,1)$\\
V & $1$ & $2$ & $0$ & $0$ & $(0,0,1,0)$\\
VI & $1/2$ & $3/2$ & $0$ & $1/4$ & $(1,1,1,0)$\\
VII & $1$ & $3$ & $-1$ & $-4$ & $(0,0,1,0)$\\
VIII & $2$ & $25/8$ & $-1/8$ & $1/2$ & $(2,2,3,2)$\\
IX & $2$ & $13/4$ & $-1/4$ & $0$ & $(1,1,2,1)$\\
X & $1/2$ & $49/32$ & $-1/32$ & $1/8$ & $(2,2,1,0)$\\
XI & $1/2$ & $25/16$ & $-1/16$ & $0$ & $(1,1,1,0)$\\
\bottomrule
\end{tabular}
\end{center}
Cases II, IV--VI, IX and XI test all relevant equality surfaces,
including the simultaneous equality $a=A$, $g=a+c$.
For $b<0,0,>0$ inside the $\mu>0$ cases, one may respectively
take $a=1/2,1,2$, and $g=1/2,1,1+a/2$ for I, II, III.
For all three signs of $b$ in VII, take the same three $a$ values
and $g=g_f+1$. These choices also check that omitted sign subcases
do not alter the counts. Internal exact-arithmetic checks supplement
these displayed examples; they are not numerical dependencies or
independent verification of the module.

For reproducibility of that finite check, take
$V=\{1/4,1/2,1,2,3,4\}$ and all $(a,A,c,g)\in V^4$.
For each $(a,A,c)\in V^3$, also take
$g\in\{c,a+c,g_f\}+\{-\eta,0,\eta\}$, where
$\eta=\min\{a,A,c,1\}/64$. The union contains 2,868 distinct
positive tuples. Set $k=A/3$, $h=2A/3$ for this grid; use the
displayed $k=h=1/2$ for the eleven table examples. Exact arithmetic
in $u+v\sqrt D$, with rational $u,v,D$, checks root residuals,
Vieta relations, domain comparisons, and every positive-count
equivalence and its converse. Including the eleven table examples,
4,438 original scalar residuals vanish exactly and 217 occurrences
of $s=-1$ are explicitly excluded. All eleven cells and all equality
surfaces are represented. These finite checks found no counterexample;
they supplement rather than establish the exhaustive proof.

\paragraph{Boundary continuity.}
The ordered formulas $s_\pm$ are continuous on the parameter set
$D\ge0$, since $A>0$. At $D\downarrow0$ they both tend to $s_f$;
for $D<0$ their absence as real roots is recorded, not bridged by
a real continuation. At $\mu\to0$ with fixed $b>0$,
$(s_-,s_+)\to(0,b/A)$; with fixed $b<0$,
$(s_-,s_+)\to(b/A,0)$. With $b=0$ and $\mu\downarrow0$, both labels
tend to zero as in \eqref{eq:P02-square-root-branches}. At $g\to c$,
$(s_-,s_+)\to(-1,a/A)$; the first root is excluded from the DM domain
at the boundary.
At $b\to0$ all limits within $D\ge0$ are again given by the same
continuous formulas. The sign-specific table is therefore compatible
with every boundary, although admissible root counts can change there.

\clearpage
\paragraph{Logical completeness and attempted counterexamples.}
The origin was retained before division, and every other equilibrium
is equivalent to a nonzero root of $Q$ in $s>-1$. This proves that
none was lost or introduced. In particular, the excluded root $-1$
at $g=c$ and the zero root at $\mu=0$ cannot be counted as additional
nonzero equilibria. Positivity, discriminant sign, and root order
partition all parameter tuples as in the theorem; each implication
in the necessary-and-sufficient statements is checked in every cell.
The independent positive-half-line argument supplies a second check.
Case VIII is an explicit canonical counterexample to the rejected
claim that $a+c\le g$ rules out positive synchrony. Case IV is a
counterexample even at equality. Case X shows that two admissible
nonzero roots need not be positive; case XI shows that a repeated
domain root need not be positive. Case V prevents counting a repeated
zero as a positive equilibrium, and case II prevents admitting a
canceled denominator root. All these distinctions are essential.

\paragraph{Priority checks for the independent P02 auditor.}
Independently reconstruct the following steps.
\begin{enumerate}[label=(\arabic*),leftmargin=*]
\item Componentwise substitution with coincident offsets and uniqueness
on the maximal interval in $\Domain$.
\item The origin before division and equivalence to $Q=0$ only for
$s\ne0$, $s>-1$.
\item All discriminant identities, root formulas, and Vieta relations.
\item The smaller-root domain test at $g=c$, including exclusion of
$s=-1$ even when the synchronous rational terms cancel.
\item Every cell, count, and multiplicity in
Theorem~\ref{thm:P02-complete}, including $\mu<0$ and negative roots.
\item Both directions of every alternative in
Theorem~\ref{thm:P02-positive} and its original-parameter corollary.
\item The branch graph, $g_f$, $s_f$, and the simultaneous degeneracy,
distinguishing multiplicity of $Q$ from that of $f$.
\item All boundary limits and the quadratic coefficient in
\eqref{eq:P02-quadratic-coefficient}.
\item The exact representative tuples and original-equation residuals,
including the excluded-root and double-counting tests.
\item The scope of weaker assumptions and the absence of any stability
or bifurcation conclusion.
\end{enumerate}
Internal checking and compilation do not change the DRAFT status.
